\documentclass[10pt,a4paper]{amsart}
\usepackage[margin=19mm]{geometry}
\usepackage{amsmath,amssymb,mathtools}
\usepackage[scr=rsfs,cal=euler]{mathalfa}
\usepackage{xfrac}
\usepackage[unicode,hidelinks,bookmarks=false]{hyperref}
\newcommand{\ie}{i.e.\ }
\newcommand{\cf}{cf.\ }
\newcommand{\resp}{respectively\ }
\newcommand{\Subsection}{\subsection}
\providecommand{\nobreakdash}{\nobreak\hbox{-}}
\setlength{\emergencystretch}{2em}
\multlinegap0pt
\usepackage{amssymb}
\usepackage{mathtools}
\usepackage{booktabs}
\newtheorem{theorem}{Theorem}
\title{Strong Central 2-Trees with Tail Degrees \texorpdfstring{$\{2,3\}$}{\{2,3\}}:\\Structural Characterization and Uniqueness Criteria}
\author{Julian Allagan, Shawn Langley, Weizheng Gao, Mohamed Elbakary}
\date{Source: arXiv:2512.18378v1 (CC BY 4.0)}
\begin{document}
\maketitle

\noindent\emph{Source and licence.} Strong Central 2-Trees with Tail Degrees \texorpdfstring{$\{2,3\}$}{\{2,3\}}:\\Structural Characterization and Uniqueness Criteria. Version arXiv:2512.18378v1 (CC BY 4.0), distributed by arXiv under the Creative Commons Attribution 4.0 International licence, http://creativecommons.org/licenses/by/4.0/, as recorded in the arXiv licence metadata for this exact version and shown on the abstract page https://arxiv.org/abs/2512.18378v1. Full licence notice: \texttt{licenses/arxiv-2512.18378-CC-BY-4.0.txt}.\par\smallskip

\noindent\textit{Editorial scope.} Counted unit: the article's theorem thm:quadratic-lower (source0.tex lines 842--851). The article's complete original proof of it is delivered with it (source0.tex lines 854--932, one \texttt{proof} environment of 79 lines). No mathematics is written, repaired or re-derived: the statement and the proof are the article's own lines, in the article's own notation, and no retained line is substituted or re-flowed.  No further source-local context is needed: every cross-reference of every retained line resolves inside the two delivered spans.  Nothing was omitted from any retained span, so the package carries no omission record.  No retained line contains a citation, so the wrapper carries no bibliography and no bibliography entry.  No retained line contains a figure environment, an image-inclusion command or drawing code, so no image is delivered and the package declares no asset dependency.  Editorial formatting only, in addition: an AMS \texttt{amsart} preamble that restates the article's own declarations and macros used by the retained lines, namely the environments \texttt{theorem} on separate counters, together with the standard packages those lines need.  The retained environments print in this document as Theorem 1 in order of appearance, whereas the article prints the same items with the numbering of its own full document.  The complete native source of the article as distributed in the arXiv e-print for this exact version is bundled unchanged as \texttt{sources/191386/source0.tex}.
\begin{theorem}[Quadratic lower bound for tricentral case]
\label{thm:quadratic-lower}
Let $N_3(n)$ denote the number of non-isomorphic strong tricentral $2$-trees on $n$ vertices whose tail degrees lie in $\{2,3\}$.
There exists an absolute constant $c>0$ such that $N_3(n)\ge c\,n^2$ for all sufficiently large $n$ with $n\equiv 0\pmod 3$.
In particular,
\[
N_3(n)\;\ge\;\left\lfloor \frac{(n-15)^2}{72}\right\rfloor
\]
for every $n\ge 24$ with $n\equiv 0\pmod 3$.
\end{theorem}

\begin{proof}
Fix an integer $K\ge 5$ and put $n=9+3K$, so $n\ge 24$ and $n\equiv 0\pmod 3$.
For integers $p,q$ with $1\le p<q\le \lfloor K/2\rfloor$, we construct a strong 
tricentral $2$-tree $G(p,q)$ on $n$ vertices, and we show that $G(p,q)\cong G(p',q')$ 
holds if and only if $(p,q)=(p',q')$.

Let $a,b,c$ span a triangle, which will be the core. On each core edge we stack 
two \emph{pages}: on $ab$ add $x_{ab}^a,x_{ab}^b$, on $bc$ add $x_{bc}^b,x_{bc}^c$, 
and on $ca$ add $x_{ca}^c,x_{ca}^a$; each page is adjacent precisely to the 
endpoints of its supporting edge. This produces a $2$-tree on $9$ vertices in 
which $\deg(a)=\deg(b)=\deg(c)=6$.

Given $(p,q)$, we now attach chains by successive stacking (each step stacks a 
new degree-$2$ vertex on an existing edge, hence preserves the $2$-tree property).
From $a$ we attach a chain of length $p$ starting at edge $\{a, x_{ab}^a\}$ and 
a chain of length $K-p$ starting at edge $\{a, x_{ca}^a\}$. From $b$ we attach 
a chain of length $q$ starting at edge $\{b, x_{bc}^b\}$ and a chain of length 
$K-q$ starting at edge $\{b, x_{ab}^b\}$. From $c$ we attach a single chain of 
length $K$ starting at edge $\{c, x_{ca}^c\}$, leaving edge $\{c, x_{bc}^c\}$ unused.

Exactly $p+(K-p)+q+(K-q)+K=3K$ vertices are added, so $|V(G(p,q))|=9+3K=n$.
Each core vertex gains exactly $K$ new neighbors, hence $\deg(a)=\deg(b)=\deg(c)=6+K$.
Every chain vertex has degree $2$ (if terminal) or $3$ (if internal), and each 
page has degree $2$ or $3$ since it supports at most one chain. Therefore all 
tail vertices have degree in $\{2,3\}$, and $G(p,q)$ is strong tricentral.

Among the core vertices, $c$ is uniquely characterized as the one for which the 
two incident page edges support chain-length multiset $\{0,K\}$ (one unused, one 
of length $K$). The other two core vertices have both incident page edges used, 
since $1\le p,q\le \lfloor K/2\rfloor <K$. Hence every isomorphism fixes $c$.
After fixing $c$, the only remaining symmetry is the transposition of $a$ and $b$.
Consequently, the unordered pair of multisets
\[
\mathcal{I}(p,q):=\bigl\{\{p,K-p\},\{q,K-q\}\bigr\}
\]
is an isomorphism invariant of $G(p,q)$.

Set $m=\lfloor K/2\rfloor$. For $1\le p<q\le m$ we have $p\le m\le K/2$, so 
$p<K-p$ and similarly $q<K-q$; in particular, $\{p,K-p\}$ and $\{q,K-q\}$ 
each consist of two distinct integers. Moreover, $p\neq q$ implies 
$\{p,K-p\}\neq \{q,K-q\}$. Thus $\mathcal{I}(p,q)$ uniquely determines the 
ordered pair $(p,q)$ up to swapping $a$ and $b$, and the constraint $p<q$ 
removes this ambiguity. Therefore $G(p,q)\cong G(p',q')$ if and only if $(p,q)=(p',q')$.

It follows that
\[
N_3(n)\ge\binom{m}{2}=\binom{\lfloor K/2\rfloor}{2}.
\]

We now establish the claimed bound by considering two cases.

Since $n=9+3K$, we have $(n-15)^2/72=(3K-6)^2/72=(K-2)^2/8$.

\textbf{Case 1:} $K=2t$ for some integer $t\ge 3$.
Then $m=\lfloor K/2\rfloor=t$ and
\[
\binom{m}{2}=\frac{t(t-1)}{2}\ge \frac{(t-1)^2}{2}
\quad\text{(using $t\ge t-1$)}
\]
\[
=\frac{(2t-2)^2}{8}=\frac{(K-2)^2}{8}
\ge \left\lfloor \frac{(K-2)^2}{8}\right\rfloor.
\]

\textbf{Case 2:} $K=2t+1$ for some integer $t\ge 2$.
Then $m=t$ and
\[
\frac{(K-2)^2}{8}=\frac{(2t-1)^2}{8}=\frac{4t^2-4t+1}{8}=\frac{t(t-1)}{2}+\frac{1}{8},
\]
so $\lfloor (K-2)^2/8\rfloor=t(t-1)/2=\binom{m}{2}$.

In both cases,
\[
N_3(n)\ge\binom{\lfloor K/2\rfloor}{2}\ge\left\lfloor \frac{(K-2)^2}{8}\right\rfloor
=\left\lfloor \frac{(n-15)^2}{72}\right\rfloor.
\]

The quadratic-growth statement follows by taking any $c<1/72$ and restricting to $n$ large enough so that $(n-15)^2/72\ge cn^2$.
\end{proof}

\end{document}
